\chapter{Introduction}
\label{ch:intro}

\section{Docking as a generative problem}

Molecular docking asks where a small molecule --- the \emph{ligand} --- sits
when bound to a protein, and in what orientation and internal conformation.
The answer determines whether a candidate compound is worth synthesising, and
computational prediction has been a central tool of structure-based drug
design for decades. Classical docking programs pose this as optimisation:
define a scoring function that approximates binding free energy, then search
the space of poses for its minimum. The weakness is well known --- the scoring
function is a crude surrogate, and the search returns a single point estimate
with no calibrated notion of how much to trust it.

Generative modelling reframes the task. Instead of minimising a surrogate
energy, one learns to \emph{sample} from the conditional distribution
$p(\text{pose} \mid \text{protein}, \text{ligand})$ implied by experimentally
determined structures. Drawing $K$ samples then yields a set of hypotheses
whose spread is itself informative. DiffDock \citep{corso2023diffdock}
established that this works at scale, and a line of subsequent work has
refined the state space, the architecture and the noising process.

What makes the problem geometrically interesting is that a pose is not a
vector. Placing a rigid body in space requires a rotation and a translation:
an element of the special Euclidean group $\SEthree$. Treating a ligand as a
small number of rigid fragments makes the configuration space a product
$(\SEthree)^N$, a compact-times-Euclidean manifold on which the usual
machinery of Gaussian noise and linear interpolation does not directly apply.
Generative models on such spaces must be built with the geometry, not
retrofitted to it.

\section{Research question and scope}

SigmaDock \citep{prat2026sigmadock} is a recent docking model that combines an
$\SEthree$-equivariant EquiformerV2 backbone \citep{liao2024equiformerv2} with
a score-based diffusion process over rigid-fragment poses. It is the starting
point of this thesis. Diffusion is not the only way to learn a transport
between a tractable source distribution and the data; flow matching
\citep{lipman2023flow} learns a velocity field directly by regression, without
simulating a stochastic process at training time, and its Riemannian extension
\citep{chen2024riemannian} carries the construction to manifolds. The two
approaches are closely related --- diffusion admits a deterministic
probability-flow ODE \citep{song2021sde} --- but they differ in what the
network is asked to predict, in how the source distribution is chosen, and in
how many network evaluations sampling requires.

This thesis asks:

\begin{quote}
\itshape
What changes when SigmaDock's score-based diffusion is replaced by Riemannian
flow matching on the protein--ligand pose manifold, while architecture, data
and computational budget are held fixed --- and which flow-specific design
choices are required before the substitution becomes competitive?
\end{quote}

The design principle throughout is \emph{minimal intervention}. The backbone,
featurisation, data pipeline, optimiser and evaluation protocol are inherited
unchanged; only the generative mechanism is replaced. This is what makes the
comparison interpretable, and Chapter~\ref{ch:sigmaflow} documents the extent
to which it was achieved, including where it was not.

Two limitations are stated at the outset rather than buried. First, all
experiments in this thesis are run at a fraction of the compute used to
produce SigmaDock's published numbers; absolute performance is therefore well
below the literature and no benchmark claim is made. What is claimed is a
\emph{matched-budget} comparison between two generative engines. Second, the
ligand is treated as a set of rigid fragments, so torsional flexibility enters
only through the fragmentation and not as a modelled degree of freedom. Both
constraints are inherited from SigmaDock and are shared by both arms of the
comparison.

\section{Contributions}
\label{sec:contributions}

\begin{enumerate}[leftmargin=*,itemsep=2pt]
\item A controlled replacement of the generative mechanism in a
      state-of-the-art docking system by Riemannian flow matching, with the
      network backbone verified byte-identical between arms
      (\S\ref{sec:minimal-change}). The construction itself applies known
      theory; the \emph{control} is what makes the comparison usable.
\item Identification and correction of two defects in the rotational
      machinery: a frame inconsistency in the predicted vector field
      (\S\ref{sec:frame-defect}) and a clamping error in the logarithmic map
      inherited from the parent codebase (\S\ref{sec:log-defect}). Both are
      documented with the behavioural signature of the fix.
\item A derivation showing that the geodesic conditional target on $\SOthree$
      is independent of time, and the resulting shrinkage argument for why
      this removes the implicit curriculum a diffusion noise schedule provides
      (Chapter~\ref{ch:asymmetry}). This is the principal theoretical
      contribution.
\item An empirical programme testing that mechanism against training compute,
      together with the exclusion of three competing explanations ---
      integration resolution, molecular symmetry and loss weighting
      (\S\ref{sec:excluded}).
\item A leakage-controlled construction for an informative source
      distribution on $\SOthree$, motivated directly by the mechanism rather
      than imported from elsewhere (\S\ref{sec:source-results}).
\end{enumerate}

\dnote{UPDATE AFTER FINAL RUNS}{Re-rank this list once
Chapters~\ref{ch:results} and \ref{ch:discussion} are complete. If the
informative-source experiment succeeds, item 5 moves up; if the rotational
deficit persists under longer training, item 3 becomes the headline.}
